\documentclass[11pt]{article}
\usepackage{ap-lit}
\title{The Exam, the Rubric, and the Three Essays}
\author{Paul Wang}
\date{November 2026}
\docpart{Part 0 \textperiodcentered\ Study Guide}

\begin{document}
\maketitle

\begin{abstract}
This part is the map for the rest of the series. It covers the format of the exam (section 1), the seven skill categories the multiple choice is written to (2), a method for reading a passage and its question set (3), the six-point analytic rubric used for all three essays (4), and how to earn each row of it: the thesis (5), evidence and commentary (6), and sophistication (7). The last section (8) is a plan for the two-hour essay section, including how to prepare a work for the Literary Argument. Each section ends with check questions; the key is on the last page. The unit guides then go genre by genre: prose fiction (Unit 1), poetry (Unit 2), and longer fiction and drama (Unit 3).
\end{abstract}

\scoringline{Check questions: \blank\ / \qtotal \qquad Target: 12+ \qquad Weakest rubric row: \blank[5em]}

\section{The exam at a glance}
\label{sec:1}

The exam takes 3 hours and is fully digital (Bluebook): you read the passages and type the essays on screen.

\begin{compare}{Section}{What it is}{Weight and time}
I. Multiple choice & 55 questions in 5 sets of 8--13, each set after one passage. At least 2 passages are prose fiction (drama can count as prose) and at least 2 are poetry. Every question has \textbf{four} choices (A--D) since the 2025 exam. & 45\% of the score; 60 minutes (about 65 seconds per question, \emph{including} reading time). \\
II. Free response & Three essays, each scored 0--6 on the same analytic rubric: \textbf{Q1 Poetry Analysis} (a given poem), \textbf{Q2 Prose Fiction Analysis} (a given passage), \textbf{Q3 Literary Argument} (a work you choose). & 55\% of the score; 2 hours for all three. Pacing is yours; about 40 minutes per essay is the usual plan. \\
\end{compare}

The course is organized in nine units that spiral through three genres. On the multiple choice, \textbf{short fiction} is 42--49\%, \textbf{poetry} 36--45\%, and \textbf{longer fiction or drama} 15--18\%. Since the free response gives you one poem, one prose passage, and one open question, both genres matter in both sections.

\begin{trap}
There is no penalty for a wrong answer. Never leave a multiple-choice question blank; if the clock is running out, guess on every remaining item.
\end{trap}

\begin{q}{3}{0.1 Exam format}{55 MCQ in 60 min, 45\%; three essays in 2 h, 55\%.}
Which statement about the exam is accurate?
\ch{The essays and the multiple choice are each worth 50\% of the score.}{Each multiple-choice question has five answer choices.}{The free-response section is worth more than the multiple-choice section.}{The Literary Argument essay is written about a passage printed on the exam.}
\end{q}

\begin{q}{1}{0.1 Exam format}{Five sets, 8--13 questions each, at least two prose fiction and two poetry passages.}
How is the multiple-choice section built?
\ch{Five passages, each followed by 8--13 questions, with at least two prose fiction and at least two poetry passages}{Ten short passages with five or six questions each}{Three long passages: one poem, one story, and one play}{Separate sections for vocabulary, grammar, and reading}
\end{q}

\section{The seven skill categories}
\label{sec:2}

Every multiple-choice question is written to one skill, and every skill asks for the same move: explain the \emph{function} of a choice (what it does in the text), not just its name. College Board's weights for the multiple-choice section:

\begin{compare}{Skill category}{What the question asks you to explain}{Weight}
1 Character & What a character's words, actions, and contrasts reveal about values, motives, and change. & 16--20\% \\
2 Setting & How time, place, and atmosphere shape character, mood, and conflict. & 3--6\% \\
3 Plot and structure & How the order of events, shifts, pacing, and form create meaning. & 16--20\% \\
4 Narrator or speaker & Who tells or speaks, from what distance, with what attitude and reliability. & 21--26\% \\
5 Word choice, imagery, symbols & How specific words, images, and symbols create tone and meaning. & 10--13\% \\
6 Comparison & How similes, metaphors, personification, and allusions transfer qualities from one thing to another. & 10--13\% \\
7 Literary argument & What claim about the text is best supported, or what evidence best supports a claim. & 10--13\% \\
\end{compare}

Narrator or speaker is the single largest category. On every passage, settle two questions before the first item: \emph{who is telling this, and how do they feel about it?} About a fifth of the questions will ask some version of that.

\begin{keyidea}[Function, not label]
``The passage uses a simile'' is a label. ``The simile comparing the boy's body to a harp presents his feeling as something played upon him from outside, not something he controls'' is a function. Questions and essays both reward the second.
\end{keyidea}

\begin{q}{4}{0.2 Skill categories}{Narrator or speaker is the largest category, 21--26\%.}
Which skill category carries the most weight on the multiple choice?
\chh{Setting}{Character}{Comparison}{Narrator or speaker}
\end{q}

\begin{q}{2}{0.2 Function}{A function names what the choice does to meaning, not just its name.}
Which sentence explains the \emph{function} of a detail rather than only identifying it?
\ch{The poem contains alliteration in line 3.}{The repeated hard \emph{c} sounds in line 3 make the speaker's anger audible just as he claims to be calm.}{The story is told in the first person.}{The passage includes several metaphors and a simile.}
\end{q}

\section{Reading a passage and its questions}
\label{sec:3}

\subsection{Before the questions}
\begin{enumerate}[leftmargin=1.6em]
\item \textbf{Read the italic introduction.} It gives author, date, genre, and sometimes the situation. A date tells you which conventions to expect (a Victorian narrator who comments freely; a modernist narrator who stays inside one mind).
\item \textbf{Read the whole passage once for the situation and the speaker.} Who is speaking, to whom, about what, and in what mood?
\item \textbf{Find the shift.} Most passages and poems turn once: a ``but,'' a new stanza, a change in tense, a move from description to reflection. Many questions cluster around it.
\end{enumerate}

\subsection{Common question stems}
\begin{terms}
``In context, the phrase \ldots\ most nearly suggests'' & Meaning in \emph{this} sentence, not the dictionary meaning. Reread the full sentence before choosing. \\
``The \ldots\ primarily serves to'' & Function. Ask what would be lost if the detail were cut. \\
``The narrator's attitude toward \ldots\ is best described as'' & Tone. Choose the pair of words that fits the whole passage, including any mixed feeling. \\
``The shift in line 9 / paragraph 4 is from \ldots\ to'' & Structure. Both halves of the answer must be right. \\
``Which of the following best supports the claim that \ldots'' & Literary argument. Match the claim's exact wording to the evidence. \\
\end{terms}

\subsection{Choosing}
Predict an answer in your own words before you read the choices. Then eliminate: a choice is wrong if \emph{any} part of it is wrong. Watch for four distractor types: \textbf{half right} (one accurate half, one inaccurate half); \textbf{too extreme} (``despair'' for mild regret, ``contempt'' for teasing); \textbf{true but not responsive} (accurate about the passage but not what the stem asks); and \textbf{too literal} (a figurative line read as fact).

\begin{trap}
Tone answers on this exam are almost never one-dimensional. When two choices both seem partly right, the answer is often the one that holds two feelings at once (``affectionate but ironic,'' ``admiring yet uneasy'').
\end{trap}

\begin{q}{1}{0.3 Distractors}{Any wrong part makes the whole choice wrong: a half-right distractor.}
A choice describes a narrator as ``nostalgic and bitter.'' The narrator is clearly nostalgic but shows no bitterness anywhere in the passage. The choice is
\ch{wrong, because one half of it is not supported}{right, because it is partly supported}{right if no other choice mentions nostalgia}{wrong only if the passage is a poem}
\end{q}

\begin{q}{3}{0.3 Reading method}{Shifts mark the turn in a passage, and many questions cluster around them.}
On a first reading, the most useful thing to locate, besides the speaker and situation, is
\ch{every unfamiliar word}{the number of sentences in each paragraph}{the point where the passage shifts in tone, time, or focus}{every example of alliteration}
\end{q}

\section{The six-point essay rubric}
\label{sec:4}

All three essays are scored with the same analytic rubric. Each row is scored separately.

\begin{compare}{Row}{What earns the points}{What does not}
A. Thesis (0--1) & A \textbf{defensible interpretation} that responds to the prompt. It may be more than one sentence if the sentences are together, and it can appear anywhere, though the introduction is safest. & Restating the prompt; a summary of what happens; listing devices with no claim about meaning; a claim that does not answer the prompt. \\
B. Evidence and commentary (0--4) & \textbf{4}: specific evidence supports all claims in a line of reasoning, and the commentary \emph{consistently} explains how the evidence supports it. On Q1 and Q2, the commentary explains how \emph{multiple} literary elements or techniques contribute to the interpretation. \textbf{3}: the same, but the commentary explains only \emph{some} of the evidence. \textbf{2}: some specific evidence, but commentary that is uneven, repetitive, or partly faulty. \textbf{1}: mostly summary or general evidence with little commentary. \textbf{0}: restates the thesis or is irrelevant. & Quotations dropped in with no explanation; paraphrase in place of analysis; naming devices without saying what they do. \\
C. Sophistication (0--1) & Sophistication of thought or a complex literary argument, shown \emph{across} the essay: exploring tensions or complexities in the text; situating the interpretation in a broader context; accounting for an alternative interpretation; or a consistently vivid, persuasive style. & A single phrase or reference; sweeping generalizations (``Since the dawn of time\ldots''); a lone sentence about ``society'' tacked onto the end. \\
\end{compare}

Two consequences follow. First, Row B is two-thirds of each essay, and the difference between a 2 and a 4 is commentary, not more quotations. Second, sophistication is rarely earned by adding something at the end; it comes from a thesis that already contains a tension (``X, yet Y'') and a body that develops it.

\begin{q}{2}{0.4 Rubric}{Row B carries 4 of the 6 points.}
Of the six points available on each essay, how many come from evidence and commentary?
\chh{1}{4}{3}{2}
\end{q}

\begin{q}{4}{0.4 Rubric}{Quotations with little explanation = mostly evidence, little commentary: Row B 1--2.}
An essay quotes eight well-chosen lines but follows each quotation with a sentence of plot summary. It most likely scores low in Row B because it
\ch{uses too many quotations}{lacks a thesis}{does not discuss the author's life}{supplies evidence without explaining how it supports the argument}
\end{q}

\begin{q}{3}{0.4 Sophistication}{Row C is earned across the argument, e.g., by exploring a tension, not by one phrase.}
Which is most likely to earn the sophistication point?
\ch{Opening with a quotation from a famous critic}{Ending with a sentence about how the theme applies to everyone}{Arguing throughout that the speaker's praise of her rival is both sincere and quietly self-serving}{Using several advanced vocabulary words}
\end{q}

\section{Row A: the thesis}
\label{sec:5}

A thesis that earns the point makes a claim about \emph{meaning} that a reasonable reader could dispute and that the text can support. A useful shape for Q1 and Q2:

\begin{keyidea}[Thesis template]
In [text], [author] uses [two or three specific choices] to portray [subject] as [complex interpretation], suggesting [insight about the work or human experience].
\end{keyidea}

Compare three openings for an essay on a story about a widow who feels relief at her husband's death:
\begin{itemize}
\item \emph{``The author uses imagery, diction, and irony to convey the widow's feelings.''} Restates the prompt; no interpretation. No point.
\item \emph{``The widow feels both grief and relief.''} An interpretation, but thin: it names the complexity without saying anything about it. Usually earns the point; leaves little to argue.
\item \emph{``Through the open window's images of spring and the narrator's careful refusal to name the widow's joy outright, the author presents a freedom so forbidden that the widow can admit it only to herself, and only in a whisper.''} Defensible, specific, and complex; it sets up the body paragraphs.
\end{itemize}

\begin{trap}
A thesis can be built around devices, but it must not \emph{be} a list of devices. The devices are the evidence; the interpretation is the claim.
\end{trap}

\begin{q}{1}{0.5 Thesis}{Only this choice makes a defensible claim about meaning.}
Which sentence would earn the thesis point for an essay on how a poem conveys a speaker's attitude toward aging?
\ch{The speaker's images of autumn and dying fire cast aging as a loss that, by sharpening a listener's love, becomes a kind of gift.}{The poem is about getting older.}{The poet uses imagery, metaphor, and couplets in the poem.}{Aging is something everyone experiences.}
\end{q}

\begin{q}{3}{0.5 Thesis}{The thesis may be more than one sentence and may appear anywhere, if it is a defensible interpretation.}
Which is true of the thesis under the rubric?
\ch{It must be exactly one sentence.}{It must be the last sentence of the first paragraph.}{It may be more than one sentence and can appear anywhere, as long as it presents a defensible interpretation.}{It earns the point if it names at least three devices.}
\end{q}

\section{Row B: evidence and commentary}
\label{sec:6}

\subsection{Evidence}
Use short, embedded quotations, a few words to a line, grammatically woven into your sentences. You are not graded on the length of a quotation; you are graded on what you do with it. On Q3, specific references to moments in the book (a scene, a decision, a final image) serve as evidence; you do not need exact quotations.

\subsection{Commentary}
Commentary answers three questions about each piece of evidence:
\begin{terms}
What? & What does the evidence say or show, briefly? \\
How? & How does the technique work: what does the word connote, what does the comparison transfer, what does the shift contrast? \\
Why? & Why does it matter to the thesis: how does it build the interpretation? \\
\end{terms}

\subsection{A worked example}
Emily Dickinson's ``Hope is the thing with feathers'' (published 1891) calls hope a bird that ``perches in the soul'' and ``sings the tune without the words.''
\begin{study}{Summary (Row B 1)}
Dickinson says hope is like a bird that sits in the soul and sings a song without words. This shows that hope is important.
\end{study}
\begin{study}{Commentary (Row B 3--4)}
By making hope a bird that merely ``perches,'' Dickinson gives it the lightness of a creature that could fly off at any moment, yet it stays. Its song ``without the words'' places hope beneath language: it cannot be argued with because it makes no argument. Hope, in this image, is fragile in body but stubborn in presence.
\end{study}

\subsection{Line of reasoning}
Organize body paragraphs by \emph{claims}, not by devices. A paragraph called ``imagery'' tends to list; a paragraph that argues ``the setting first imprisons her'' and the next that argues ``the same window becomes her way out'' builds a line of reasoning that the rubric rewards.

\begin{q}{2}{0.6 Commentary}{Commentary explains how the wording works and why it matters.}
Which sentence is the strongest commentary on the phrase ``a summons to all my foolish blood''?
\ch{The narrator uses the word ``summons.''}{A ``summons'' is an order one cannot refuse, so the boy presents his love as an authority over him, while ``foolish'' shows the older narrator judging the feeling he still describes.}{This quotation shows that the boy likes the girl.}{The phrase is an example of figurative language.}
\end{q}

\begin{q}{4}{0.6 Line of reasoning}{Claims, not devices, should organize paragraphs.}
Which plan for body paragraphs best builds a line of reasoning?
\ch{Paragraph 1: imagery. Paragraph 2: diction. Paragraph 3: tone.}{Paragraph 1: the beginning. Paragraph 2: the middle. Paragraph 3: the end.}{One paragraph with every quotation in the order it appears}{Paragraph 1: the house first confines her. Paragraph 2: the same house becomes the stage for her defiance.}
\end{q}

\section{Row C: sophistication}
\label{sec:7}

The rubric names four routes. You need one, developed through the essay.
\begin{enumerate}[leftmargin=1.6em]
\item \textbf{Complexity or tension}: the text holds two things at once (love and resentment; comedy and cruelty), and your essay explains how.
\item \textbf{Broader context}: you place the reading in a larger frame (genre conventions, the period, another reading of the work) in a way that sharpens it rather than decorating it.
\item \textbf{Alternative interpretations}: you acknowledge a plausible reading and show why yours accounts for more of the text.
\item \textbf{Vivid, persuasive style}: precise and consistently controlled prose. This is the hardest route to rely on alone.
\end{enumerate}

\begin{q}{1}{0.7 Sophistication}{Acknowledging and answering a plausible counter-reading is a named route.}
A student argues that a story's ending is hopeful, then notes that it could be read as bleak and explains why the closing image supports the hopeful reading more fully. This mainly shows
\ch{accounting for an alternative interpretation}{plot summary}{a thesis that restates the prompt}{evidence without commentary}
\end{q}

\section{Planning the free-response section}
\label{sec:8}

\subsection{Q1 and Q2 (given texts)}
Spend about 8--10 minutes reading and planning, then about 30 writing. Read the prompt first, so you read the text for the thing you must explain. Mark the shift, two or three strong word choices, and the ending. Write the thesis before the body, then a short introduction, two or three body paragraphs of claim--evidence--commentary, and a brief conclusion that extends rather than repeats.

\subsection{Q3 (Literary Argument)}
The prompt names an idea (a character torn between duties, an act of deception, a journey home) and asks how it contributes to an interpretation of the work \emph{as a whole}. Prepare two or three works you know closely: their central conflicts, three or four key scenes, the ending, a symbol or two, and the main characters' changes. Then choose the one that fits the prompt best, not the one you like most. Your essay must reach the meaning of the whole work; plot summary earns little.

\begin{trap}
On Q3, the most common failure is a paragraph that retells a scene. Every reference to the plot should be followed by an explanation of what it shows about the prompt's idea and the work's meaning.
\end{trap}

\begin{q}{2}{0.8 Q3 preparation}{Q3 rewards depth in a well-known work, connected to the work as a whole.}
The best preparation for the Literary Argument is to
\ch{memorize long quotations from a dozen novels}{know two or three substantial works closely, including key scenes, conflicts, and endings}{read plot summaries of many works the night before}{plan to write about whatever work the exam prints}
\end{q}

\begin{q}{3}{0.8 Planning}{Reading the prompt first focuses the reading on what must be explained.}
On Q1 and Q2, why read the prompt before the text?
\ch{The prompt includes the answers to the multiple choice.}{The prompt lists the devices that must appear in the essay.}{The prompt tells you which complexity to read for, so your annotations serve the thesis.}{The prompt is worth points of its own.}
\end{q}

\makekey
\end{document}
